\section{Computational study}\label{sec:computations}

The experiments address four questions: whether the exported cuts are
valid in practice; whether they change SCIP's results on benchmark
models; how often the targeted structure occurs and what the separator
costs; and how much of SCIP's root gap the cuts close on instances that
have the structure of Section~\ref{sec:composition}. We report three
campaigns on MINLPLib models and one on a constructed family. The first
two campaigns were run with time budgets fixed in advance; the third
repeats the comparison with the final implementation, longer budgets,
several seeds and a structure-selected sample. All three are reported in
full, including failures.

\subsection{Setup}\label{sec:setup}

All runs use SCIP~10.0.2 through PySCIPOpt~6.2.1, Python~3.13, one SCIP
thread and one BLAS thread per process, on a 36-core Linux host that was
shared with unrelated jobs (one-minute load average between about 10 and
17). Timings are therefore descriptive; we compare modes only within
the same job, where the modes run back to back. Models come from a cached
copy of MINLPLib \citep{BussieckDrudMeeraus2003} in OSiL format, restricted
by size (Appendix~\ref{app:instances}) and stratified by the library's
metadata into convex models, nonconvex models without integer variables,
and nonconvex models with integer variables.

A model counts as \emph{solved} if SCIP reports optimality or the relative
gap limit $10^{-4}$, and the returned incumbent passes an independent check
of all original rows, bounds, domains, integrality restrictions and the
objective at a scaled tolerance of $10^{-5}$; the checker evaluates the
source expressions without the model builder. Times are compared by
shifted geometric means with shift one second over models solved in every
mode. Dual bounds are compared at a relative tolerance of $10^{-4}$, the
gap limit; smaller differences are within the noise of a single run (we
report the effect of the tolerance below). Every recorded cut is replayed
in a fresh process against the archived source model and the stored SCIP
row, and deliberately corrupted records check that replay rejects
altered inputs. A run whose worker failed has no complete cut log; its cut
count is reported as unknown, not as zero.

\subsection{Campaign 1: lifted cuts with auxiliary variables}\label{sec:campaign1}

The first implementation added an auxiliary variable $w=f(x)$ for each
selected expression of one or two variables and separated support cuts in
$(x,w)$, using the polygon and star oracles, Bernstein bounds and ball
arithmetic. Because the auxiliary variables change the model, it was
compared with a \emph{control} mode that adds them without cuts. The
campaign used 24 MINLPLib models selected by a fixed hash ordering, 13
constructed mechanism cases, a six-second budget and a two-second root
budget.

Of the 24 models, 20 passed the model import in every mode; three were
rejected because a domain restriction could not be proved on the declared
box, and one failed with an unsupported variable power. Baseline and
control solved 19 of the 20, both cut modes 18; the lost model,
\code{genpooling\_lee2}, reached the time limit in both cut modes. On the
models solved in every mode, the shifted geometric mean time was
$0.38$\,s for the baseline, $0.52$\,s for the control, and $0.57$\,s and
$0.46$\,s for the two cut modes. On the mechanism cases, the control and
both cut modes solved all 13 and the baseline 12, so the additional solve
is explained by the reformulation, not by the cuts. All 1{,}082 recorded
cuts passed replay, and all 270 returned incumbents passed the
original-model check. The auxiliary variables thus cost time on their own,
and the cuts did not recover it. This motivated the original-variable cuts
of Section~\ref{sec:original}.

\subsection{Campaign 2: cuts in the original variables}\label{sec:campaign2}

The second implementation is the one described in
Section~\ref{sec:implementation}. Its prospective campaign used 30 new
MINLPLib models, ten per stratum, selected by a fixed hash ordering from
the 422 eligible models that had not been used before, with a 30-second
budget, root runs with one node, 13 mechanism cases, ten diagnostic models
from the first campaign, and seed-one repeats on six models (282 runs).

All 30 models were admitted, and every mode solved the same 25. The five
unsolved models received no cut in any mode. Cuts were added on 7 models
in mode \code{all} (38 cuts) and on 3 in mode \code{auto} (10 cuts), all of
which the baseline also solved; four of them were solved at the root in
every mode. At the gap-limit tolerance, no final or root dual bound was
better in a cut mode than in the baseline; four time-limited models had
worse final bounds in both cut modes. At a tolerance of $10^{-6}$ two models
appeared better and up to seven worse, but these differences are of the
same size as the variation between identical baseline reruns. The cut
modes were slower: on the 25 commonly solved models the shifted geometric
mean time rose from $0.55$\,s to $0.99$\,s (\code{all}) and $0.96$\,s
(\code{auto}), and both cut modes were more than 10\% slower on 21 of the
25. Discovery accounted for 24.4 of the 27.6 callback seconds of mode
\code{all}; direction LPs took 0.7\,s and certification 2.4\,s.

The campaign also exposed two defects of the frozen code. On two
diagnostic models, symbolic polynomial conversion over all model symbols
crashed during discovery (four worker errors, with unknown cut logs), and
discovery did not check the time allowance between rows, so it exceeded
the allowance in 46 cut-mode runs, by up to 19\,s. Both were repaired, and
a matched cohort of the 25 affected model and seed groups, 75 runs, was
rerun with all three modes. It had no worker errors; every mode solved the
same models; and the largest remaining overrun of the separator allowance
was 0.03\,s. On \code{graphpart\_clique-40}, for example, the worse cut-mode
bound of the original runs disappeared after the repair. All 123 cuts of
the prospective campaign and all 42 of the repair cohort passed replay,
and all 338 returned incumbents passed the original-model check.

\subsection{Campaign 3: corrected code, longer runs, structured models}\label{sec:campaign3}

Campaign~3 uses the repaired implementation with its frozen limits and
three parts fixed in a protocol before any run. Part~A repeats the 30
models of campaign~2 with a 300-second budget and seeds 0, 1 and 2, plus
root runs with one node. Part~B uses a sample selected by structure: from
the 392 eligible models not used in campaign~2, the 111 on which discovery
finds a block that the \code{auto} rule admits, ranked by a hash of a
fixed prefix and the name; we take the first 30, seeds 0 and~1. Part~C is
the constructed family of Section~\ref{sec:mechanism}. The root runs of
Parts~A and~B include a diagnostic mode \code{all-diag} with raised limits
(128 blocks, 200 cuts, 50 per callback, 10 callbacks, 500 support calls,
30\,s), which tests whether the frozen limits, rather than the cuts, bound
their effect. Parts~A and~B together comprise 690 runs.

\begin{table}[t]
\centering\small
\caption{Campaign 3 on MINLPLib models: runs solved, recorded cuts (runs
with cuts), and shifted geometric mean time over the runs solved by every
mode. Part~A: the 30 models of campaign~2, seeds 0--2, 300\,s. Part~B: 30
structure-selected models, seeds 0--1, 300\,s.}\label{tab:campaign3}
\begin{tabular}{@{}llrrrr@{}}
\toprule
Part & Mode & Solved & Cuts (runs) & SGM time (s) & Runs faster than baseline\\
\midrule
A & baseline & 80/90 & 0 & 1.05 & ---\\
  & all      & 80/90 & 151 (33) & 1.27 & 11/80\\
  & auto     & 80/90 & 37 (15) & 1.21 & 6/80\\
\midrule
B & baseline & 52/60 & 0 & 1.33 & ---\\
  & all      & 52/60 & 179 (40) & 1.87 & 8/52\\
  & auto     & 52/60 & 145 (38) & 1.82 & 9/52\\
\bottomrule
\end{tabular}
\end{table}

Table~\ref{tab:campaign3} summarizes the full runs. In both parts and in
every seed, all modes solved the same models. In Part~A, the longer budget
solved \code{graphpart\_clique-40}, which had timed out in campaign~2, and
\code{tln7} in two of three seeds; \code{ex5\_2\_5}, \code{ex8\_3\_4} and
\code{waterx} remained unsolved in every mode and seed. Cuts were added on
11 of the 30 models in mode \code{all} and on 5 in mode \code{auto}. At the
gap-limit tolerance the final dual bounds of the 90 paired runs were
better in 2 and worse in 4 for \code{all}, and better in 3 and worse in 3
for \code{auto}; between seeds, the baseline itself differed in 3 of 30
comparisons, so these differences are within run-to-run variation. The cut
modes were slower by 15--21\% in shifted geometric mean and by a median
factor of 1.3 per run.

In Part~B, all 30 models were admitted, and in each seed the same 26 were
solved in every mode; \code{bayes2\_20}, \code{bayes2\_30}, \code{bayes2\_50}
and \code{kall\_circles\_c8a} timed out. Cuts were added on 20 of the 30
models. The final bounds were better in 2 of 60 comparisons in each cut
mode, both on \code{kall\_circles\_c8a} at the time limit, and never worse.
At the root the cuts had a visible effect on four models: they closed the
root gap of \code{pooling\_bental4tp} in every cut mode and that of
\code{pooling\_bental4pq} in modes \code{auto} and \code{all-diag}, which
were then solved at the root node, and
they improved the root bounds of \code{ex3\_1\_4} (from $-6$ to $-5.79$,
and to $-5.69$ with raised limits; the optimum is $-4$) and of
\code{pointpack04}. In two root comparisons, \code{auto} ended with a worse
root bound (\code{pooling\_haverly2pq} and \code{pointpack04}): valid cuts
cannot weaken an LP, but they change SCIP's subsequent separation and its
cut selection, and the final root bound need not improve. These effects did
not change any full solve, because SCIP solves these models in about a
second regardless, and the cut modes were about 40\% slower in shifted
geometric mean and slower on 44 of 52 runs. Raising the limits at the root
(\code{all-diag}) added more cuts, 403 in Part~A and 358 in Part~B, but
improved the root bound beyond the gap tolerance on only one more model in
each part (\code{cvxnonsep\_psig20r} in Part~A and \code{pooling\_bental4pq}
in Part~B).

All returned incumbents of campaign~3, 679 in Parts~A and~B and 80 in
Part~C, passed the original-model check. Replay passed for all recorded
cuts: 7{,}498 cuts in Parts~A to~C (188 in the full and 465 in the root runs of
Part~A, 845 in Part~B and 6{,}000 in Part~C), and in each part all 14
corrupted records were rejected.

\subsection{How common is the structure, and what does it cost?}\label{sec:coverage}

The structural scan of Part~B ran discovery alone on the 392 eligible
models outside the campaign-2 sample. The importer admitted 389 of them
(three timed out while building the model). Discovery found at least one
admissible block in 234 models, a quadratic block of dimension at least two
in 185, a block with at least two nonlinear rows in 110, and a block that
the \code{auto} rule admits in 111; none of the latter are convex models.
Discovery took a median of 0.04\,s and at most 4.2\,s per model. The
structure is therefore common in MINLPLib; what is rare is a model in which
it is decisive for SCIP.

The separator's cost is dominated by Python-level work. In campaign~2,
before the repair, discovery accounted for 24.4 of 27.6 callback seconds in
mode \code{all}. In campaign~3, discovery is fast and exact certification
dominates: in Part~B, certification took 23.0 of 32.9 callback seconds in
mode \code{all}, the direction LPs 3.9\,s and discovery 4.4\,s. Over the 60
full runs this is about half a second per run, which is of the same order as
SCIP's whole solving time on most of these models.

\subsection{A family with the path structure}\label{sec:mechanism}

The path family of Section~\ref{sec:interleave} is a natural test of
whether the cuts supply strength that SCIP lacks. We built 20 instances,
five for each $n\in\{10,20,40,80\}$, each consisting of $n$ independent copies
of~\eqref{eq:family} with $\omega_A=\omega_C=1$, coupled by one nonbinding
row $\sum_iy_i\le0.8n$: variables $x_i,y_i,z_i\in[0,1]$ and $t_i$, rows
$D_i(x_i,y_i,z_i)-t_i\le0$, objective $\min\sum_it_i$. For each copy, four
distinct values are drawn from $\{k/64:k=0,\dots,48\}$ and split into two
interleaving two-point sets with random orientation. By
Theorem~\ref{thm:interleave} the optimum is $\sum_i\delta_i^2/2$, which we
know exactly. The protocol, fixed before any run, compares the baseline
with mode \code{all} under limits proportional to $n$ ($n$ blocks, $4n$ cuts,
$n$ per callback, 10 callbacks, $20n$ support calls, 60\,s), in root runs
and in full runs of 300\,s.

\begin{table}[t]
\centering\small
\caption{The path family: median fraction of the root gap closed,
$(\text{root bound with cuts}-\text{baseline root bound})/(\text{optimum}-\text{baseline root bound})$,
and instances solved within 300\,s (of five per $n$). ``Frozen'' is the
separator as used in all campaigns; ``row direction'' is the post hoc
variant described in the text, with the same limits and with wider
limits.}\label{tab:mechanism}
\begin{tabular}{@{}rccccccc@{}}
\toprule
& \multicolumn{3}{c}{Root gap closed (median)} & \multicolumn{4}{c}{Solved in 300\,s}\\
\cmidrule(lr){2-4}\cmidrule(l){5-8}
$n$ & frozen & row dir. & row dir., wide & baseline & frozen & row dir. & row dir., wide\\
\midrule
10 & 0.35 & 0.44 & 1.00 & 5 & 5 & 5 & 5\\
20 & 0.25 & 0.44 & 1.00 & 4 & 4 & 5 & 5\\
40 & 0.19 & 0.43 & 1.00 & 0 & 0 & 0 & 5\\
80 & 0.22 & 0.46 & 1.00 & 0 & 0 & 0 & 5\\
\bottomrule
\end{tabular}
\end{table}

SCIP's root bound on these instances is weak: between $-0.031$ and $-0.022$ per copy on average, while
each copy contributes at least $1/8192$ to the optimum, and SCIP closes the
gap by branching on the leaves, which its presolve turns into binaries
(Section~\ref{sec:path-example}). The frozen separator improved the root
bound on all 20 instances but closed only a median of 19--35\% of the root
gap (Table~\ref{tab:mechanism}). In full runs both modes solved the same 9
instances; on these the cuts reduced the number of nodes, for example from
30{,}798 to 2{,}316 and from 166{,}693 to 69{,}996, and the shifted geometric
mean time from 10.3\,s to 6.7\,s. All other instances timed out in both
modes.

The reason for the partial closure is in the direction search, not in the
support computation. For each block, the exact support in the direction of
its own row, $a$ equal to the row's coefficients of the block variables and
$\lambda=1$, is the block minimum, and the $n$ resulting cuts
$t_i\ge\min D_i$ make the root bound equal to the optimum. The frozen
separator's first directions, one per row, set $a=0$ and therefore bound
only the nonlinear remainder of the row; the row's affine terms in the block
variables enter the cut only through the elimination. In addition, the
first blocks of a round take up to four cuts each, so with $n$ cuts per
callback later blocks receive none. After seeing the results of Part~C, we
ran a diagnostic, reported separately from the prospective results: the
same separator with the row direction tried first, once with the same
limits and once with $4n$ cuts per callback, $16n$ cuts and $40n$ support
calls. The row direction alone raised the median closure to 43--46\%; with
the wider limits, the root gap closed completely on all 20 instances, and
all 20 full runs were solved, mostly at the root node, in 3 to 39\,s, against
9 of 20 for a rerun of the baseline. The diagnostic also ran the row
direction on the root runs of Parts~A and~B: the root bounds of all 60 models were unchanged at the gap tolerance;
the variant changed the number of cuts in 8 of 180 cut-mode runs and no
bound. The conclusions for the benchmark models therefore do not depend on
this defect.

This family was constructed to have the structure of
Section~\ref{sec:composition}, and its objective is a sum of block
functions, which is the most favorable case for block support cuts. It
shows that exact block support can supply strength that SCIP's root
relaxation lacks, that the strength reaches the solver only if the
direction search tries the right directions, and that, on such instances,
it can turn time-outs into root solves. It is not evidence about other
models.

\subsection{Summary and limitations}\label{sec:comp-summary}

Across three campaigns and the constructed family, every recorded cut
passed replay against the source model and the row stored by SCIP, and
every returned incumbent passed the original-model check. On benchmark
models the cuts never changed which models SCIP solved, in any campaign,
mode or seed. They improved root bounds on a few pooling and small
nonconvex models that SCIP solves in a second anyway, and they cost time:
in the final implementation about 15--20\% in shifted geometric mean on the
campaign-2 models and about 40\% on the structure-selected models. On the
constructed family they turned time-outs into root solves once the
direction search tried each row's own direction. We therefore do not
recommend enabling the separator by default.

The study has limitations that the reader should weigh. The test sets are
small models (at most 120 variables), because the exact oracles are
limited to small blocks and the implementation is in Python; larger models
may contain more or less useful structure. Time limits are short by the
standards of MINLP benchmarking, although 300\,s is long relative to the
solve times of most of the models used. The host was shared, so timing
differences of a few percent are not meaningful; we have reported them
only as shifted geometric means over paired runs, with the variation
between baseline seeds for comparison \citep{LodiTramontani2013}. Only SCIP
was used, because the separator is a SCIP component; the baseline is SCIP
with all its default nonlinear machinery, which is the relevant comparison
for such a component. Finally, the post hoc diagnostic of
Section~\ref{sec:mechanism} was designed after the prospective results were
seen and is reported as such; it explains the mechanism results and does
not change the benchmark conclusions.
